% !TeX root = ../technical_report.tex
\chapter{Repository Architecture, Directory Hierarchy, and File Extension Semantics}
\label{ch:repository_architecture}

The engineering of a high-throughput, epistemically safe Machine Learning Operations (MLOps) platform requires an uncompromising foundational architecture. This chapter enumerates the physical layout, the manifest invariants, and the toolchain standardisation that govern the TaxonVision-MLOps repository. By enforcing strict boundary conditions between immutable source code, declarative infrastructure manifests, and volatile biological datasets, the architecture guarantees deterministic reproducibility across all development, continuous integration, and production environments. The repository must be approached not merely as a collection of source files, but as a rigorous state machine governed by continuous integration gates and mathematically verified file extension semantics.

\section{Architectural Philosophy and Clean Separation}
\label{sec:architectural_philosophy}

The TaxonVision-MLOps repository is designed as a unified monorepo. However, it deliberately eschews the chaotic entanglement frequently observed in monolithic data science repositories. Instead, the architecture enforces a rigorous, physical separation of concerns through its directory topology. This separation is not merely a stylistic convention; it is a fundamental security and operational invariant that dictates how Continuous Integration (CI) pipelines interact with the underlying code. The core philosophy dictates that stateless domain logic, stateful configuration, containerised deployment specifications, and mutable data artifacts must never intersect within the same directory tree. This decoupling ensures that a modification to an inference parameter does not trigger a recompilation of the application image, and that an update to the taxonomic catalogue does not invalidate the mathematical tests of the core domain logic.

\subsection{Isolation of Domain Logic}
The functional core of the application resides exclusively within the \texttt{src/taxon\_vision/} directory. This directory is strictly reserved for pure, stateless Python source code. It encapsulates the domain models representing biodiversity entities, the mathematical formulations for loss functions, the network inference graphs, and the application programming interface routing logic. By isolating the computational logic, the system ensures that the core application remains agnostic to its deployment environment and its underlying data sources.

\subsection{Externalisation of State and Configuration}
Configuration and state are deliberately externalised from the core logic. The \texttt{config/} directory serves as the definitive source of truth for all declarative state, including hierarchical parameters, the taxonomic catalogue, and Prometheus telemetry definitions. This separation guarantees that parameters dictating model behaviour can be adjusted without modifying the underlying application logic. The decoupling of parameters allows continuous experimentation without source code mutation.

\subsection{Decoupling of Execution Infrastructure}
The deployment infrastructure and execution automation are similarly segregated. The \texttt{deploy/} directory contains the Kubernetes manifests required to instantiate the service within a distributed cluster, completely decoupled from the application logic itself. Furthermore, the \texttt{scripts/} directory houses the operational runbooks and continuous integration automation scripts, ensuring that the mechanisms for testing, validating, and deploying the software are maintained independently of the software being deployed. This directory includes essential tools for enforcing code invariants, such as \texttt{validate\_okf.py} for Open Knowledge Format conformance, and \texttt{audit\_license\_compliance.py} to ensure biological imagery adheres strictly to open-access legal mandates.

\subsection{Separation of Mutable Data Artifacts}
Perhaps the most critical separation is the exclusion of large, mutable binary artifacts from the version control system. The \texttt{data/} directory serves as a lightweight manifestation of the remote data topography. Rather than storing physical images or extracted embeddings, it houses cryptographic pointer files governed by Data Version Control (DVC), such as \texttt{sample.dvc}. This isolation protects the Git repository from unbounded bloat whilst simultaneously ensuring bit-for-bit reproducibility of the datasets used during training and validation. The \texttt{.dvc/config} ensures that pointers securely map to the remote DagsHub S3 blob storage, maintaining the lineage of the data without incurring the heavy storage cost within the Git tree itself.

\section{Directory Anatomy and Manifest Invariants}
\label{sec:directory_anatomy}

The physical layout of the repository is formalised in Table~\ref{tab:directory_anatomy}. Each directory acts as a bounded context, governed by specific ownership rules and pre-commit security invariants.

\begin{table}[htbp]
\centering
\caption{TaxonVision-MLOps Directory Anatomy and Security Invariants}
\label{tab:directory_anatomy}
\begin{tabularx}{\textwidth}{@{}l l X@{}}
\toprule
\textbf{Directory Path} & \textbf{Primary Function} & \textbf{Security Invariant / Ownership} \\
\midrule
\texttt{src/taxon\_vision/} & Core domain logic, models, and API. & Strict MyPy type checking; No mutable state allowed. \\
\texttt{config/} & Declarative configuration parameters. & Defines systemic state. \\
\texttt{data/} & Ephemeral data caching and DVC pointers. & Strictly excluded from Git via \texttt{.gitignore}; DVC tracks MD5 hashes. \\
\texttt{deploy/} & Deployment configuration and infrastructure. & Formalised declarative definitions for zero-downtime cluster deployment. \\
\texttt{docs/} & OKF v0.2 Knowledge Base and LaTeX reports. & Enforced Open Knowledge Format validations and mathematical notation formatting. \\
\texttt{scripts/} & CI/CD automation and local runbooks. & Shell and Python utility scripts; Isolated from core module imports. \\
\texttt{tests/} & Pytest suites and conformal invariants. & Minimum coverage thresholds enforced; Executes in ephemeral sandboxes. \\
\texttt{.agents/} & Autonomous AI agent routing and roles. & Governed by \texttt{index.md} algorithmic routing matrices and role assignments. \\
\bottomrule
\end{tabularx}
\end{table}

\subsection{Root Configuration Manifests}
The root of the repository houses the definitive, authoritative manifests that control the execution environment, the automated task sequences, and the code quality gates. These manifests are strictly enforced; any deviation or manual override within a local environment is considered a violation of the system's operational integrity.

\textbf{The Pixi Project Specification (\texttt{pyproject.toml} and \texttt{pixi.lock}):} The environment is governed by a unified \texttt{pyproject.toml} manifest, which integrates the declarative Pixi configuration under the \texttt{[tool.pixi.*]} namespace. This specification formally defines dual environments: a \texttt{dev} environment featuring GPU acceleration via CUDA 12.x for local training, and a \texttt{ci-dev} environment targeting CPU-only execution for continuous integration workers. This dual-environment specification ensures that the dependency graph remains cryptographically identical across hardware platforms, whilst selectively provisioning heavy GPU libraries only where required. The deterministic locking of this environment is handled by the \texttt{pixi.lock} file.

\textbf{The Declarative Recipe Runner (\texttt{Justfile}):} Operational automation is orchestrated via a \texttt{Justfile}. This file replaces disparate shell scripts with formal, parameterised recipes. Commands abstract the complexity of the underlying toolchain, guaranteeing that every developer and automated agent executes operations using identical, validated command sequences.

\textbf{The Quality Assurance Gate (\texttt{.pre-commit-config.yaml}):} The repository enforces an aggressive pre-commit pipeline containing twenty distinct automated quality hooks. These hooks encompass trailing whitespace elimination, secret detection via Gitleaks, strict MyPy type assertion, and bespoke internal validations such as \texttt{verify\_markdown\_rules.py} for formatting compliance, \texttt{validate\_okf.py} for checking the Open Knowledge Format, and \texttt{audit\_license\_compliance.py}. This automated gating mechanism ensures that no non-compliant code, malformed documentation, or restrictive biological imagery licenses can mathematically enter the source tree.

\section{File Extension Semantics and Format Specifications}
\label{sec:file_extension_semantics}

To maintain high throughput and strict reproducibility, the architecture mandates specific file extensions and standardises upon high-performance serialisation formats. Each extension carries formal semantic weight and processing invariants.

\subsection{Data and Model Serialisation}

\textbf{\texttt{.dvc} (Data Version Control Pointers):} Large binary datasets, such as raw citizen-science imagery and extracted feature matrices, are physically tracked outside the Git tree. Instead, Git tracks lightweight \texttt{.dvc} pointer files (e.g., \texttt{sample.dvc}). These declarative manifests contain the MD5 cryptographic hash, the byte size, and the relative path of the corresponding remote artifact without bloating the Git tree. When an environment is synchronised, the DVC client resolves these hashes against the DagsHub S3 remote and materialises the assets into the local \texttt{data/} directory. This mechanism guarantees exact dataset reproducibility without imposing bloat upon the version control repository.

\textbf{\texttt{.parquet} (Columnar Observation Tables):} Tabular metadata, such as the iNaturalist taxonomic hierarchies and observation records (\texttt{observations.parquet}), are exclusively stored in the \texttt{.parquet} format. Apache Parquet is a columnar serialisation format that employs dictionary encoding and Snappy compression. Unlike traditional row-based formats (such as CSV or relational database exports), Parquet allows the data processing engines to perform projection pushdowns, scanning only the required taxonomic columns. This results in massive reductions in I/O latency and memory allocation when aggregating millions of species occurrences.

\textbf{\texttt{.onnx} (Open Neural Network Exchange):} Deep learning computational graphs are serialised and deployed using the \texttt{.onnx} Protocol Buffers format. The architecture exports trained architectures to ONNX, which provides a hardware-agnostic, zero-allocation computational graph. The ONNX format allows the integration of dynamic execution providers and supports aggressive INT8 quantisation, reducing the memory footprint by up to 75\% whilst accelerating vector multiplication routines.

\textbf{\texttt{.pt} (PyTorch Serialised States):} During the active learning and training phases, intermediate network states and fine-tuned taxonomic classifier heads are persisted as \texttt{.pt} files. These artifacts contain the ordered dictionaries of the network weights and the momentum states of the AdamW optimiser. However, these files are strictly ephemeral; they are tracked via MLflow or DVC for lineage and provenance, but they are never deployed directly to the inference serving infrastructure.

\subsection{Configuration and Architecture Definitions}

\textbf{\texttt{.toml} (Declarative Configuration):} The repository standardises on the \texttt{.toml} format (Tom's Obvious, Minimal Language) for dependency and project configuration. This format provides an unambiguous, hierarchical syntax that avoids the parsing complexities of YAML and the rigidity of JSON. The central \texttt{pyproject.toml} manifest specifies the entire Python application state, the Ruff linting rules, the Pytest coverage configurations, and the Pixi environment resolutions. Furthermore, the \texttt{zensical.toml} defines the documentation generation parameters.

\textbf{\texttt{.yaml} (Hierarchical Manifests):} While \texttt{.toml} handles project dependencies, \texttt{.yaml} is exclusively reserved for structural configuration. This includes the Hydra-compatible configuration matrices in the \texttt{config/} directory, the declarative Kubernetes workloads, the Jenkins continuous integration pipeline definitions, and the \texttt{.pre-commit-config.yaml} hook declarations.

\textbf{\texttt{.lock} (Cryptographic Dependency Sealing):} The dependency graph is deterministically sealed within the \texttt{pixi.lock} file. This file contains the precise, cryptographically hashed versions of every upstream Conda and PyPI package across all required operating systems and hardware platforms (Linux-64 CPU and GPU). This locking mechanism eliminates "works on my machine" syndromes and guarantees that the containerised inference service executes with exactly the same dependency tree as the developer workstation.

\textbf{\texttt{.tikz} (Vectorised Architectural Schematics):} The documentation system leverages the \texttt{.tikz} extension to define standalone vector graphic descriptions. Rather than relying on external, opaque image editors, the architectural schematics are authored in raw LaTeX TikZ code. This programmatic approach ensures that all diagrams remain infinitely scalable, visually consistent with the surrounding typography, and fully trackable through version control.

\section{Toolchain Interlock and Environmental Isolation}
\label{sec:toolchain_interlock}

The operational integrity of the TaxonVision-MLOps pipeline is heavily dependent upon the standardisation of its execution toolchain. To this end, the architecture completely abandons traditional Python dependency management tools.

\subsection{The Prohibition of Standard Dependency Managers}
The repository strictly enforces a ban on the direct invocation of \texttt{pip}, \texttt{poetry}, \texttt{uv}, and the creation of standard \texttt{.venv} directories. These traditional tools, while prevalent, fail to manage non-Python dependencies (such as the underlying CUDA libraries, OpenBLAS, or system-level compression tools) in a unified, deterministic manner.

\subsection{Pixi Environment Isolation}
Instead, the entire environment is constructed and isolated using Pixi. Pixi builds upon the Conda ecosystem, providing a fully cross-platform, multi-environment specification mechanism. When a developer or CI pipeline executes \texttt{just setup}, Pixi resolves the \texttt{pyproject.toml} specification and builds an isolated dependency tree within the \texttt{.pixi/} directory. This environment encompasses not only Python packages but also the precise versions of the NVIDIA CUDA Toolkit, the ONNX Runtime execution providers, and foundational binary libraries.

\subsection{Deterministic Execution and the Frozen Graph}
All commands must be executed within the context of this environment. By doing so, the system guarantees that the Python interpreter, the mathematical libraries, and the deep learning frameworks are sourced exclusively from the sealed \texttt{pixi.lock} graph. This toolchain interlock is mathematically verifiable; any attempt to circumvent the Pixi wrapper or to install ad-hoc dependencies will result in a failure of the continuous integration tests, thereby protecting the epistemic safety of the ultimate production deployment.

\section{Extended Architectural Rigour and Process Constraints}
\label{sec:extended_rigour}

To further guarantee the robustness of the system, the architectural philosophy extends beyond mere directory structure and file formats. It imposes rigorous process constraints on how state transitions occur within the MLOps lifecycle. This formalisation ensures that the evolution of the model, from raw data ingestion to production deployment, follows a verifiable and mathematically sound trajectory.

\subsection{Immutability of the CI/CD Pipeline}
The Continuous Integration and Continuous Deployment (CI/CD) pipeline acts as the ultimate arbiter of code quality. Local developer environments are considered inherently untrustworthy. Therefore, any pull request targeting the main branch must successfully traverse the entire suite of automated checks. This includes not only standard linting, but also the execution of necessary invariant tests. Failure to satisfy these mathematical invariants results in an immediate pipeline failure, preventing compromised logic from reaching the deployment stage.

\subsection{The Role of the Model Registry}
The MLflow Model Registry serves as the authoritative ledger for model lineage. When a training pipeline completes successfully, the resulting classification head is not blindly deployed. Instead, it is registered within MLflow alongside its complete metadata profile. This profile includes the Git commit hash of the executing code, the DVC hash of the dataset used for training, the precise hyperparameters selected, and a comprehensive suite of evaluation metrics. This rigorous lineage tracking ensures that every prediction made in production can be traced back to its exact origin, providing complete auditability for ecological researchers.

\subsection{Agentic Governance and Roles}
The repository governance extends to autonomous AI agents through the \texttt{.agents/} directory. The foundational \texttt{AGENTS.md} serves as the root index, whilst the routing matrix in \texttt{.agents/index.md} enforce a rigid protocol for automated codebase modifications. Agents are restricted by path-scoped roles, ensuring they only modify files within their designated domain. This agentic governance structure guarantees that automated refactoring or documentation updates do not violate the architectural invariants established for the core modules.

\subsection{Deployment Automation and Zero-Downtime Tolerances}
The final stage of the pipeline involves deploying the trained and quantised ONNX models to the Kubernetes cluster. The manifests within the \texttt{deploy/} directory formalise this process. The deployment configuration leverages declarative methodologies to ensure zero-downtime transitions between model versions. Readiness and liveness probes continuously monitor the health of the inference service. If a newly deployed pod fails to initialise the ONNX runtime or fails to respond to health checks within the specified timeout, the deployment is automatically rolled back. This declarative approach to infrastructure management eliminates the risk of human error during deployment operations and guarantees the continuous availability of the species identification service.

\subsection{Conclusion on Architectural Integrity}
The culmination of these architectural constraints -- the strict separation of concerns, the declarative manifest definitions, the specific file extension semantics, the uncompromising toolchain interlock, and the automated quality gates -- results in a system of unparalleled robustness. The TaxonVision-MLOps repository is not merely a collection of scripts; it is a mathematically verifiable, self-contained ecosystem designed to process complex biological data streams with absolute deterministic reliability. This uncompromising approach to software engineering is the prerequisite for achieving the high levels of epistemic safety and operational maturity demanded by modern global biodiversity monitoring platforms. By strictly adhering to these principles, the architecture provides a solid and immutable foundation upon which the subsequent layers of data processing, model training, and active learning can be safely constructed.
